\section{Trabalhos Relacionados}

A literatura sobre vulnerabilidades e priorização de correções pode ser organizada em três linhas principais: predição de exploração, crítica ao uso isolado de métricas de severidade e integração de dados heterogêneos. Trabalhos iniciais demonstraram que atributos públicos de vulnerabilidades podem apoiar modelos de classificação e predição de exploração, superando heurísticas simples em determinados cenários \cite{bozorgi2010_beyond_heuristics}. Entretanto, estudos baseados na NVD também indicaram limitações no uso de uma única fonte para prever riscos ou vulnerabilidades futuras, reforçando a necessidade de integrar informações adicionais ao registro técnico das CVEs \cite{zhang2011_nvd_predict,zhang2015_nvd_cyber_risks}.

Outra linha relevante discute a diferença entre severidade técnica e exploração observada. Allodi e Massacci mostram que métricas como CVSS, embora úteis para estimar gravidade, não devem ser interpretadas como proxy direto de exploração ou risco operacional \cite{allodi2012_preliminary,allodi2014_case_control}. De forma complementar, Sabottke et al. exploram sinais públicos externos aos registros tradicionais, evidenciando que informações contextuais podem melhorar a predição de exploração real, ainda que apresentem desafios de volatilidade e reprodutibilidade \cite{sabottke2015_social_media}. A formalização do EPSS consolidou essa perspectiva ao propor uma métrica probabilística voltada à estimativa de exploração, complementando métricas de severidade como o CVSS \cite{jacobs2019_epss,jacobs2020_remediation,jacobs2021_epss}.

Além da predição, estudos recentes destacam problemas de qualidade, inconsistência e integração entre bases públicas de vulnerabilidades. Jiang et al. mostram que divergências entre repositórios podem afetar decisões de priorização e remediação \cite{jiang2021_inconsistency}, enquanto abordagens baseadas em grafos de conhecimento indicam que relações entre CVE, CWE, CPE e outras entidades podem apoiar análises mais estruturadas de risco \cite{shi2024_threat_kg,yin2024_vulkg}. Em paralelo, trabalhos voltados a ecossistemas open source reforçam a importância de representar vulnerabilidades no nível de pacotes, versões afetadas e versões corrigidas, especialmente em bases como OSV e GitHub Advisory Database \cite{wu2024_affected_libraries,ruohonen2025_osv_propagation,mandl2026_osv_groundtruth}.

Diante desse cenário, este trabalho posiciona-se como uma proposta de integração reprodutível de dados públicos de vulnerabilidades, combinando fontes voltadas à identificação, severidade, exploração conhecida e advisories de pacotes. Diferentemente de abordagens centradas em uma única base, o dataset proposto busca preservar a proveniência dos dados, registrar metadados de coleta e preparar uma estrutura extensível para análises futuras de exploração, risco e priorização de remediação.

\begin{table}[htbp]
    \centering
    \caption{Síntese dos trabalhos relacionados e lacunas abordadas}
    \label{tab:related_work}
    \small
    \renewcommand{\arraystretch}{1.1}
    \resizebox{\textwidth}{!}{%
        \begin{tabular}{|>{\RaggedRight\arraybackslash}p{3.6cm}|>{\RaggedRight\arraybackslash}p{4.6cm}|>{\RaggedRight\arraybackslash}p{4.8cm}|>{\RaggedRight\arraybackslash}p{4.0cm}|}
            \hline
            \textbf{Referência}                                                                 & \textbf{Contribuição} & \textbf{Limitação observada} & \textbf{Relação com este trabalho} \\
            \hline
            Bozorgi et al.; Zhang et al.                                                        &
            Uso de dados públicos para predição de exploração e risco.                          &
            Uso limitado de fontes isoladas, principalmente NVD.                                &
            Motiva a construção de uma base multi-fonte.                                                                                                                                    \\
            \hline
            Allodi e Massacci; Sabottke et al.                                                  &
            Avaliação da relação entre severidade, sinais públicos e exploração real.           &
            CVSS isolado não representa exploração; sinais externos podem ser voláteis.         &
            Justifica combinar severidade com sinais de exploração e contexto.                                                                                                              \\
            \hline
            Jacobs et al.                                                                       &
            Proposição e formalização do EPSS para estimar probabilidade de exploração.         &
            EPSS mede probabilidade de exploração, não risco completo.                          &
            Motiva enriquecimento futuro com score probabilístico.                                                                                                                          \\
            \hline
            Jiang et al.; Nguyen e Massacci                                                     &
            Análise de inconsistências e limitações de qualidade em bases públicas.             &
            Metadados podem ser incompletos, divergentes ou desatualizados.                     &
            Reforça a necessidade de proveniência e controle de qualidade.                                                                                                                  \\
            \hline
            Shi et al.; Yin et al.                                                              &
            Uso de grafos de conhecimento para relacionar CVE, CWE, CPE e entidades correlatas. &
            Nem sempre integram sinais operacionais como KEV, EPSS e advisories de pacotes.     &
            Apoia a modelagem futura do dataset como grafo de vulnerabilidades.                                                                                                             \\
            \hline
            Wu et al.; Ruohonen e Ramadan; Mandl et al.                                         &
            Análise de vulnerabilidades em ecossistemas open source, pacotes e versões.         &
            Escopo frequentemente centrado em ecossistemas específicos.                         &
            Motiva integração com bases orientadas a pacotes, como GHSA e OSV.                                                                                                              \\
            \hline
        \end{tabular}%
    }
\end{table}
